\scanpage{184}
\margin{\(\begin{array}{lll}\operatorname{Tr}x=b&\operatorname{Tr}y=b'&\operatorname{Tr}z=b''\\\det x=c&\det y=c'&\det z=c''\\\operatorname{Tr}yz=t&\operatorname{Tr}zx=t'&\operatorname{Tr}xy=t''\end{array}\)}
la formule
\[
\tau+\tau'=\operatorname{Tr}(xy+yx)z\qquad\text{où}
\]
\[
xy+yx=\underbrace{\operatorname{Tr}(y)}_{b'}x+\underbrace{\operatorname{Tr}(x)}_b y+\bigl(\underbrace{\operatorname{Tr}(xy)}_{t''}.\unsure{1}-\underbrace{\operatorname{Tr}x}_b\underbrace{\operatorname{Tr}y}_{b'}\bigr)1
\]
donc
\[
\deleted{\unsure{(xy)z+}}\quad
(xy+yx)z=b'xz+byz+(t''-bb')z\qquad\text{d'où}
\]
\[
\boxed{\tau+\tau'=\unsure{bt|b''t''}+b't'+b''t''-bb'b''}
\overset{\mathrm{def}}{=}\Sigma(b,b',b'',t,t',t'')
\]
On devine que $\tau\tau'$ sera également un polynôme en $b,b',b'',c,c',c'',t,t',t''$, de sorte que $\tau,\tau'$ s'expliciteront (à permutation près) comme racines d'un polynôme \inserted{unitaire} du second degré, dont les coeff $B$ ($=\tau+\tau'$) et $C$ ($=\tau\tau'$) sont des polynômes en les variables $\uline{b^{(i)}},\uline{c^{(i)}},\uline{t^{(i)}}$.

\deleted{Mais plutôt que de chercher une expression pour $\tau\tau'$, je vais essayer de les \inserted{($\tau$ et $\tau'$)} déterminer directement en termes de $\det,1$ \uline{et de $e$}.
Je suis guidé par le fait que $xy$ et $yx$ sont opposés mod $1,x,y$ --- donc \uline{si $z$ est orth.} à $1,x,y$ \deleted{\unsure{det}} \inserted{(pour $\varphi$)} on aura
\[
\operatorname{Tr}(xyz)=-\operatorname{Tr}(yxz).
\]
Pour se ramener à ce cas, on construit $\alpha,\beta,\gamma\in k$ tels que
\[
z'=z-\alpha1-\beta x-\gamma y
\]
soit orth. à $1,x,y$, i.e.
\[
\begin{aligned}
b''-2\alpha-b\beta-b'\gamma&=0\\
t'-\unsure{b}\alpha+(2c-b^2)\beta-t''\gamma&=0\\
t-\deleted{\unsure{b\alpha}}\inserted{b'\alpha}-t''\beta+(2c'-b'^2)\gamma&=0
\end{aligned}
\]
i.e
\[
\left\{\begin{aligned}
2\alpha+b\beta+b'\gamma&=b''\\
b\alpha+(b^2-2c)\beta+t''\gamma&=t'\\
b'\alpha+t''\beta+(b'^2-2c')\gamma&=t
\end{aligned}\right.
\]
matrice
\diagram{\(\begin{pmatrix}2&b&b'\\b&b^2-2c&t''\\b'&t''&b'^2-2c'\end{pmatrix}\)}
}
\margin{183}
